\chapter{Conclusion and Future Work}
\label{ch:conclusion}

\section{Summary}

This dissertation develops \systemname, a modular prototype for cross-media person identification. The system uses visual detection to obtain person crops, VLM prompting to extract structured attributes, dense and sparse vectorization to represent those attributes, and a hybrid matcher to assign identity decisions. Recent engineering work reorganized the codebase into a maintainable package structure and added a GUI-based video testset creation workflow.

\section{Contributions}

The current dissertation draft contributes:

\begin{enumerate}
  \item a working VLM-assisted person identification pipeline;
  \item a hybrid dense--sparse representation for attribute-grounded matching;
  \item a dynamic registry for open-ended attribute dimensions;
  \item a GUI workflow for converting videos into labelled crop datasets;
  \item an Overleaf-compatible paper-based dissertation structure for continuing the writing process.
\end{enumerate}

\section{Future Work}

The immediate next step is to generate and annotate larger test datasets using the new GUI. After that, the dissertation should add final experiments, including dense-only and sparse-only ablations, threshold tuning, text-query retrieval, and failure-case analysis. A future system version may integrate face embeddings, stronger tracking, caching, and local VLM inference to reduce latency and improve reproducibility.

\section{Closing Remarks}

The project began as an implementation-heavy prototype and is now ready to become a dissertation argument. The central argument is that VLM-derived attributes can make person matching more interpretable and more adaptable to cross-media workflows, provided that the system is evaluated carefully and its limitations are made explicit.
